\section{Conclusion}\label{sec:conclusion}

Integer dimension has both an algebraic and a geometric aspect. For fixed
quadratic systems, noncommutative rank determines the exact logarithmic
precision coefficient. At finite accuracy, contact covariance leads to a
determinant benchmark and constructive bounds with explicit dimension
losses. Positive structure, effective input coordinates and the actual
error body can sharpen those losses, while correlated domains can prevent
naive transfers.

For scalar convex graphs, finite chord geometry makes binary linear lifts
uniformly competitive with arbitrary convex integer lifts. Certified
integration, rational endpoint algorithms and a common index compiler
turn this comparison into polynomial-size constructions for the specified
polynomial and power representations. Coupled outputs require simultaneous
control: original-output or positive-polar curvature bases give such control,
and a shared basis across separable inputs preserves the relevant packing
lower bound. These results depend on convexity and the stated error geometry.
The exact product counts, positive-power obstruction, nonconvex polynomial
family and tilted-body example identify different ways those assumptions
matter.

Several questions remain. The finite covariance benchmark has an
order-$n\log n$ gap on an explicit family; a different efficiently
computable invariant might give smaller losses, subject to the proved
count-approximation hardness. Local and global smooth Hessian ranks can
differ, and rank-changing scalar maps are not characterized by the
constant-rank theorem. For one input and arbitrarily many convex outputs
with box error, the present logarithmic-rank upper comparison and strict
one-bit example leave open whether a universal constant additive gap
exists. The cap-set restriction strengthens one lower-bound family but
does not decide that question. Extending the separable vector construction
to mixed-coordinate nonlinear terms also requires a new compact encoding
argument.

The encoding separation supplies a further boundary: even four binary
variables can require exponentially long rational linear data relative
to a binary exponent, whereas second-order-cone constraints encode the
same graph accuracy with short data. Integer count, continuous size,
coefficient length, construction time and optimization time should
therefore remain separate when these formulations are used or compared.
